\hypertarget{discussion-threats-to-validity-and-conclusion}{%
\section{Discussion, Threats to Validity, and Conclusion}\label{discussion-threats-to-validity-and-conclusion}}

\label{sec:9}

\hypertarget{interpretation}{%
\subsection{Interpretation}\label{interpretation}}

\label{sec:9.1}

The results converge on a single message: for pre-deployment criticality analysis of pub-sub
middleware, \emph{how} a component is critical is at least as important as \emph{whether} it is --- and the case
for graph learning is narrower, and more specific, than we expected when we set out. Four findings
carry this.

\textbf{First, learning pays --- most defensibly for set identification (RQ1).} The typed predictor leads
the strongest non-learning baseline on both metrics and under all three evaluation protocols. Out of
distribution --- the genuine pre-deployment condition --- it reaches \(\rho = 0.608\) against \(0.521\) for a
QoS-weighted centrality that requires no training, no labels and no transfer assumption, and
\(F_1@K = 0.465\) against \(0.308\). Both margins are measured against a baseline we first had to repair:
\texttt{Topo-QoS} was computing no QoS weighting whatsoever (Section~\ref{sec:8.1}), and until that was fixed it was
\texttt{Topo-BL} wearing a different label. We place substantially more weight on the second metric than the
first, for three reasons. Operationally, an architect hardens a handful of components, not a ranked
list of 150. Empirically, the ranking comparison is the less stable one: \texttt{Topo-QoS} carries the
largest across-fold variance of any predictor in the study (\(\sigma = 0.31\) LOSO, \(0.34\) k-fold),
whereas the set-identification ordering does not invert under any protocol. And evidentially, the
ranking margin is the weaker of the two: it does not survive a paired test across scenarios
in-distribution (\(p = 0.375\), Table~\ref{tab:19}), and out of distribution it rests on a run whose artifact was
not retained, against an earlier retained run that recorded a tie (Section~\ref{sec:8.1}). The practical reading is
therefore a scope condition with an asymmetry in confidence: if a team needs the critical set, typed
learning is worth its training cost on the evidence here; if they need a cheap ordering, QoS-weighted
centrality remains a serviceable default, and we cannot presently demonstrate that learning beats it.

\textbf{Second, decomposition is worth having for reasons that are not accuracy (RQ2, Section~\ref{sec:8.2}; weighting
sensitivity from RQ3's robustness analysis, Section~\ref{sec:8.3}).} The dimension
weighting does not improve ranking --- equal weights beat the calibrated ones (Section~\ref{sec:8.3}) --- and the
stratified check we ran to detect Simpson's-paradox masking did not find it in the \(Q\)--\(I\) relation.
What survives is the property we think actually motivates the decomposition: a four-dimensional
profile says \emph{why} a component is critical and routes the finding to an owner, which a scalar cannot,
and that holds regardless of how the four are combined. The methodological discipline was also not
wasted: pooled-versus-stratified reporting \emph{did} catch a real distortion elsewhere in this study,
where collapsing Application and Library nodes into one correlation moved a headline figure by 0.38
(Section~\ref{sec:5.5}). The check earned its place by catching something, just not where we pointed it.

\textbf{Third, edge criticality benefits from being measured rather than assumed.} Replacing a hand-chosen
bridge multiplier with actual edge-removal simulation reversed the intuition it encoded: most
individual links turn out to be replaceable, and a whole class of structurally non-redundant edges
(\texttt{RUNS\_ON}) carries no measurable impact at all because the cascade model cannot express their
failure. The heuristic would have assigned exactly those edges their source node's full blast radius.
This is a small result with a general moral --- a plausible label-generating assumption is not a
substitute for the observation it stands in for.

\textbf{Fourth, remediation is now verified per edit, which is a stronger guarantee than the previous
aggregate but not yet a demonstration of value (Section~\ref{sec:6.7}).} Each candidate is simulated in isolation and
admitted only if it beats the simulator's own noise at every propagation threshold, so a regressing
edit can no longer be carried by an improving one. What that filter reveals is that the operators'
yield is highly topology-dependent --- from 3 of 35 candidates on Autonomous Vehicle to 38 of 58 on IoT
Smart City --- and that the resulting risk reductions are uniformly small. We regard this as the
correct outcome of an honest test rather than a failure of the mechanism: the previous design
reported a more favourable aggregate precisely because it never asked each edit to justify itself.

\textbf{Finally, automated quality gating operationalises these checks continuously (RQ4).} By evaluating
in-memory via the \texttt{MemoryRepository} and bypassing database round-trips, the framework runs
anti-pattern scans and counterfactual simulations in seconds (\textasciitilde5 s medium, \textasciitilde40 s xlarge). That speed
makes the analyzer viable as a blocking CI/CD check, and the delta-aware gate semantics (Section~\ref{sec:6.6}) make
it sustainable: it blocks newly introduced architectural regressions at commit time --- bridging the
Architecture--Code Gap --- without repeatedly flagging known, risk-accepted structure, in the manner of
``Clean as You Code'' static-analysis gates. Of the four contributions this is the one least disturbed
by the audit, and the one we would defend most confidently.

\hypertarget{threats-to-validity}{%
\subsection{Threats to Validity}\label{threats-to-validity}}

\label{sec:9.2}

\textbf{Construct validity.} D1 and D2 define criticality as Quality-in-Use loss, and this study never
observes Quality-in-Use. The validation chain has two links, and only the first is measured:

\begin{equation}
\underbrace{\text{structural / learned score}}_{Q(v),\ \text{HGL}}
\;\xrightarrow{\ \text{\textcircled{1}}\ }\;
\underbrace{\text{simulated failure impact}}_{I^*,\ I_{\text{comp}},\ I_{\text{dyn}}}
\;\xrightarrow{\ \text{\textcircled{2}}\ }\;
\underbrace{\text{real Quality-in-Use loss}}_{\text{what D1 and D2 define}}
\end{equation}

Link \textcircled{1} is what Section~\ref{sec:8} reports: a real, falsifiable result. Link \textcircled{2} is not measured anywhere in this
paper --- no user study, expert elicitation, or production incident record is used, and the simulator
is itself a \emph{model} of stakeholder harm rather than an observation of it. The defensible claim is
therefore: \emph{RMAV and the learned predictors track simulated failure impact, and simulated failure
impact is our stated operationalisation of Quality-in-Use loss.} The stronger claim --- that these
scores track Quality-in-Use as stakeholders would report it --- is not supported by anything here, and
we do not make it. Closing link \textcircled{2} requires evidence of a different kind: expert ranking studies on
the same topologies, or post-hoc comparison against incident records from a deployed system (Section~\ref{sec:9.3}).

Two qualifications keep this from being either overstated or unduly bleak. The ISO/IEC 25019
characteristics are not equally out of reach: Effectiveness and Efficiency are in principle
measurable from quantities the simulators already produce --- delivery rate before and after a fault,
and the latency shift the discrete-event engine records --- so link \textcircled{2} is partly closable by
re-summarising existing output on the Quality-in-Use axis rather than by new instrumentation.
Freedom from risk is blocked by the corpus rather than by the method: deadline and lifespan violation
counters exist and the harness has an oracle slot for them, but no topic in the scenario corpus
declares a deadline, so the counters never fire. Acceptability and Satisfaction are behavioural and
are not measurable by these means at all, which bounds what this construct can ever claim on them.
None of that has been \emph{run}; it establishes measurability, not a measurement.

Because the ground truth is simulated rather than observed, the strongest claims we can make remain
comparative: which modelling choices perform better under identical conditions, not absolute
predictive accuracy in operation. Four further bounds apply, and we state them rather than leave them
implicit.

\emph{The two oracles agree weakly.} \(I^*(v)\) and \(I_{\text{comp}}(v)\) correlate at mean \(\rho = 0.394\)
(Section~\ref{sec:7.5}). Results established against one do not transfer to claims measured against the other, which
constrains this paper's own internal cross-referencing: Section~\ref{sec:5.4}'s library finding and Section~\ref{sec:5.5}'s stratified
check are \(I_{\text{comp}}\) results and are not evidence about the \(I^*\)-backed tables in Section~\ref{sec:8.1}.

\emph{Six instrument defects were found and corrected during this revision.} All were silent --- none
raised an exception or produced an obviously wrong number --- and all are recorded here because each
had, or could have had, a published figure resting on it. The first two predate the corpus
regeneration of Section~\ref{sec:7.1}. First, the \texttt{Topo-QoS} baseline was applying no QoS weighting: \(w(t)\) is declared
on the Topic node, the harness looked for it on the pub-sub relationship, and the generated
topologies carry none there, so every derived dependency edge kept a unit weight and the baseline
computed plain betweenness on all seven scenarios. It has been repaired to resolve \(w(t)\) from the
shared Topic; the affected columns of Tables~\ref{tab:18} and~\ref{tab:20} and the k-fold table were recomputed, and the
non-QoS variants were verified unchanged to machine precision. Second, HGT attention extraction
captured nothing, because \texttt{HGTConv} in the pinned PyTorch Geometric release exposes no
\texttt{return\_attention\_weights} argument and the extraction fell through its own error branch; attention
is now captured from the layer's own softmax, and the attention subgraph of Fig.~\ref{fig:6} is generated
from real per-edge \(\alpha\) rather than an edge-weight fallback. We note that the second defect had masked a
third --- the subgraph renderer itself raised on a \texttt{networkx} API change, which nothing had exercised
while the attention payload was empty.

Four further defects were found in a later pass that checked the implementation against this
manuscript directly, rather than against a specific reported number.

Third, and the one that changes reported figures, \texttt{FaultInjector}'s cascade iterated an unordered
Python set of subscribers while consuming seeded random draws for each; set iteration order in
Python is salted per-process by \texttt{PYTHONHASHSEED}, so the \emph{same} requested seed could assign different
draws to different subscribers across processes, making \(I^*(v)\) reproducible only within one
interpreter run, not across runs --- exactly the kind of instability the seed-mean-and-standard-deviation
protocol of Section~\ref{sec:5.1} and Section~\ref{sec:7.5} was designed to average over, not diagnose. It has been repaired (the
iteration is now sorted); cross-process reproducibility was verified directly (identical \(I^*(v)\)
across five different \texttt{PYTHONHASHSEED} values, both on the synthetic corpus and on the real-world
Autoware sweep of Section~\ref{sec:8.5}). Two figures rest on the pre-fix labels and are flagged rather than silently
carried forward: the label test--retest \(\rho\)/Jaccard ceiling of Section~\ref{sec:7.5}, restated above with the
corrected, now process-independent values (previously reported as \(\rho \in [0.928, 1.000]\), Jaccard
\(\in [0.56, 1.00]\), both measured within a single process and so blind to this defect); and Section~\ref{sec:8.5}'s
Autoware row, whose ``sweep-to-sweep instability'' was reported in an earlier draft as a property of
that graph and is corrected there to what it actually was. Tables~\ref{tab:18} and~\ref{tab:20}, and every scenario in
Table~\ref{tab:13}, are unaffected: both use \texttt{cascade\_depth\_limit=0}, the setting at which the sixth defect
below is provably a no-op, and neither exercises the code path this defect lived in independently of
that setting.

Fourth, \texttt{extract\_rmav\_scores\_dict} --- the function that turns \texttt{PredictionService}'s RMAV output into
the GNN's auxiliary training target --- keyed its lookup by an attribute (\texttt{component\_id}) that the
underlying dataclass does not have (it has \texttt{id}), so every key fell through to the object's own
\texttt{repr()} string and the lookup silently returned nothing usable; the \(0.1\)-weighted RMAV-consistency
term of Table~\ref{tab:17} was training against an all-zero target wherever this function was on the path. It
has been repaired to key by \texttt{id} first, matching its sibling function's already-correct convention.
Table~\ref{tab:3}/5/6/7's reported runs are unaffected: both evaluation harnesses (\texttt{cli/loso\_evaluate.py},
\texttt{cli/kfold\_evaluate.py}) read RMAV scores through a different loader that never called the broken
function. Any GNN checkpoint trained via the standalone \texttt{cli/train\_graph.py} entry point without an
explicit \texttt{-\/-rmav} file did go through the broken path and trained with no RMAV supervision; that
entry point is not what produced the tables in this paper.

Fifth, the parallel worker in the prescription stage's per-edit verifier (Section~\ref{sec:6.4}) constructed its own
evaluator with default settings --- layer \texttt{system}, no GNN checkpoint --- regardless of what layer and
checkpoint the run was actually configured with, so a \texttt{-\/-layer\ app} run would score every candidate
edit's counterfactual impact on the \texttt{system} layer while its baselines were measured on \texttt{app}. It has
been repaired to thread the configured layer and checkpoint into each worker. Table~\ref{tab:13} is unaffected:
\texttt{reproduce/run\_prescribe\_all.py} runs at \texttt{layer="system"} with no checkpoint, which is exactly what
the unpatched default constructed, so the mismatch could not occur for the reported run.

Sixth, the post-loop computation of \(I^*(v)\) read a \texttt{topic\_loss} variable left over from the cascade's
last executed wave rather than recomputing it against the final set of failed components, so a
subscriber failure in the final wave was not reflected in that subscriber's own reported feed loss.
It has been repaired to recompute once more after the loop terminates. This defect is a no-op when
\texttt{cascade\_depth\_limit=0} (unlimited waves, the default and the setting \texttt{reproduce/} and the corpus
generation in Section~\ref{sec:7.1} both use throughout): we verified this directly by running the pre-fix and
post-fix simulators against the same cached topologies and confirming a maximum absolute difference
in \(I^*(v)\) of exactly \(0\) across three scenarios. It is not a no-op under a finite
\texttt{cascade\_depth\_limit}, which no reported figure in this paper uses.

\emph{A third of each system is unlabelled.} The cascade model cannot express the failure of a Topic or a
physical Node, leaving 30--47\% of components per scenario without ground truth. Predictions for them
are produced but never validated. Broker labels are degenerate in three of seven scenarios for a
related reason. Any claim of coverage across ``all five component types'' would be unsupported, and
the per-type results report those strata as undefined rather than as zero.

\emph{Reported figures approach the labels' own reproducibility.} The ground truth agrees with itself at
test--retest \(\rho\) of 0.807--1.000 and top-\(K\) Jaccard of 0.44--1.00 across seeds (post the
determinism fix above; these are now stable across \texttt{PYTHONHASHSEED}, unlike the figures an earlier
draft reported). A model scoring near the former has saturated the labels rather than underperformed,
and every top-\(K\) metric inherits the latter's churn.

\emph{The behavioural oracle is delivery-based, not QoS-aware.} \(I_{\text{dyn}}\) carries the
construct-validity argument of Section~\ref{sec:7.5}, so the limits of what it measures bound that argument too. Its
discrete-event engine implements deadline, lifespan, and reliability enforcement, but resolves topic
QoS from an attribute key the generated corpus does not write, so every run in this evaluation falls
back to defaults and the deadline and best-effort drop paths are structurally zero rather than
measured as zero. Latency is likewise uninformative here: at the corpus's publication rates,
utilisation stays far below saturation and queues never build, leaving \(p95\) latency flat to within
run-to-run jitter across faulted components. \(I_{\text{dyn}}\) should therefore be read as a
\emph{throughput} oracle --- it corroborates that the cascade ranking tracks lost message delivery, and it
makes no claim about QoS contract conformance under load.

\textbf{Internal validity.} The chief internal risk is circular validation --- a predictor scoring well
because its inputs leaked from its labels. The framework addresses this by \emph{view} separation:
predictors operate on \(G_{\text{analysis}}\) while ground truth is generated by simulating
\(G_{\text{structural}}\), no simulation output is fed back as a predictor feature, and remediation
candidates are generated without reading simulated impact. \textbf{This is view independence, not independence of
data source}: both views are deterministic functions of the same input topology, so what is ruled
out is feature--label feedback, not the possibility that both encode a shared modelling assumption.
The distinction matters for how much weight the guarantee can bear, and we prefer to state it than to
let ``independent simulator'' imply more.

The behavioural oracle narrows this, and it is worth being precise about by how much. A sharper form
of the circularity objection is that \(I^*\) is an artifact of its own traversal --- that a
topology-derived score is being validated against labels manufactured by walking the same topology.
\(I_{\text{dyn}}\) answers that specific charge: it reaches its ranking by simulating message traffic
through queues over simulated time, never traversing \texttt{DEPENDS\_ON}, and it recovers \(I^*\)'s ordering
(Section~\ref{sec:7.5}). The cascade \emph{algorithm} is therefore not the artifact. What remains unaddressed is the layer
beneath it: all three oracles are simulation rather than observed failure data, and all three are
deterministic functions of the same generated topology. A modelling assumption shared by the
architecture model itself would be invisible to every one of them. Calibration against instrumented
deployments (Section~\ref{sec:9.3}) is the only thing that reaches it.

Two further internal-validity issues surfaced during a pre-submission audit of this work and are
disclosed because they invalidated previously reported numbers. First, the evaluation harness scored
different predictor families on different node populations and different samples (Section~\ref{sec:7.3}); the
correction changed the sign of the RQ1 conclusion. Second, the Leave-One-Scenario-Out sweep reused
stale model checkpoints and was therefore not training at all --- the same command produces
\(\rho = -0.576\) in 3.2 s against a dirty workspace and \(\rho = +0.594\) in 322 s against a clean one.
Both are fixed and all reported figures come from the corrected runs, but the episode is itself a
finding about this class of experiment: a silently-cached artifact is indistinguishable from a
trained one in the output, and only the implausible wall-clock time exposed it.

\emph{Artifact retention is uneven across the reported tables, and one headline table cannot currently be
regenerated.} Table~\ref{tab:18} and the sensitivity sweeps of Section~\ref{sec:8.3} regenerate exactly from stored result
files --- a claim that held only approximately before the determinism defect above was fixed, since a
re-run in a fresh process was not guaranteed to reproduce a stored \texttt{FaultInjector} label exactly
even at an unchanged seed. It now holds without qualification.
The Leave-One-Scenario-Out result file behind Table~\ref{tab:20} does not exist: it was overwritten during the
revision, and the most recent retained log for that sweep predates the baseline repair and records a
different ordering (Section~\ref{sec:8.1}). We disclose this rather than present Table~\ref{tab:20} on the same footing as
Table~\ref{tab:18}, and we regard it as the direct continuation of the two defects above. The common mechanism
in all three is that an experiment's \emph{evidence} and its \emph{output} were allowed to come apart --- a
cached checkpoint, a mismatched sample, an unretained result file --- and in each case the number
looked entirely ordinary. The discipline this study now imposes, and did not impose soon enough, is
that no figure enters the manuscript unless the artifact that produced it is retained and the figure
can be recomputed from it. Table~\ref{tab:20} is the outstanding exception, and re-running it under the final
apparatus is the first item of remaining work.

\textbf{External validity.} This is the weakest dimension of the study, and we regard it as the
highest-value follow-up (Section~\ref{sec:9.3}).

\emph{The corpus spans ten deployment domains, but only three architectures are not ours.} Seven scenarios
come from a single statistical topology generator; the three real-world graphs (Autoware.universe
ROS 2, the Cloud-Native Microservices mesh, and Train-Ticket) are transcribed from published
open-source architectures. On the latter, SaG achieves mean rank correlation over five seeds of
\(\rho = 0.688\), \(0.778\) and \(0.759\), and up to \(F_1@K = 1.000\) on two of the three.

\emph{None of the three clears the framework's own gate.} All three fail SPOF-F1; Autoware additionally
fails the \(\rho\) threshold and Cloud Microservices the predictive-gain threshold, for 5 failed checks
of 15 (Section~\ref{sec:8.5}). The \(F_1@K = 1.000\) figures are partly a tie-breaking artifact of \(K\) exceeding the
count of genuinely non-zero-impact components. We read the result as evidence that the predictive
\emph{ranking} transfers beyond the generator to independently-sourced architectures --- not as a
demonstration of production readiness.

\emph{The paradigm count is two, not three.} Train-Ticket and Cloud Microservices are both microservice
meshes, so cyber-physical pub-sub is represented by Autoware alone. Leave-One-Scenario-Out evaluation
confirms inductive transfer across held-out \emph{synthetic} architectures only, since all seven share a
generator. Expanding to further middleware paradigms and to hardware-in-the-loop deployments is
future work (Section~\ref{sec:9.3}).

\textbf{Conclusion validity.} Criticality scores and simulated impact metrics exhibit heavy-tailed,
non-parametric distributions that violate normality assumptions. To prevent classification bias, we
apply non-parametric rank correlations (Spearman \(\rho\)), top-\(K\) Jaccard metrics, and adaptive
box-plot thresholding (\(Q3 + 1.5\,\mathrm{IQR}\)) rather than parametric z-scores or arbitrary absolute cutoffs (Section~\ref{sec:4.4}).

\hypertarget{limitations-and-future-work}{%
\subsection{Limitations and Future Work}\label{limitations-and-future-work}}

\label{sec:9.3}

Several limitations point to concrete next steps, ordered here by how much they would change the
paper's claims.

\textbf{Real-world deployment validation and HIL execution.} Section 8.5 narrows the external-validity gap
on three real-world open-source software architectures (Autoware.universe ROS 2, the Cloud-Native
Microservices Mesh, and Train-Ticket), but does not close it: the graphs are hand-transcribed, their
ground truth is still simulated, and none clears the framework's own validation gate in full.
Validating predictions against runtime hardware-in-the-loop (HIL) fault injection on physical
testbeds, and against harvested rather than transcribed architectures, remains the highest-value
follow-up.

\textbf{Out-of-distribution ranking is not yet a solved problem.} Section~\ref{sec:8.1} shows that a training-free
QoS-weighted centrality matches the learned models on LOSO rank correlation, with the learned
advantage confined to critical-set identification. Whether that ceiling reflects the difficulty of
cross-architecture transfer, the noise floor of the labels (Section~\ref{sec:7.5}), or a limitation of the
architecture is not resolved by these experiments. Distinguishing those three explanations ---
plausibly by training on a substantially larger and more diverse scenario corpus --- would determine
whether typed learning has more to offer here than it currently demonstrates.

\textbf{The dimension weighting does not improve accuracy.} Section~\ref{sec:8.3} finds equal weights outperform the
calibrated AHP weighting with no plateau in the shrinkage parameter. We have repositioned RMAV as an
attribution mechanism accordingly, but a weighting \emph{derived} rather than asserted --- fitted to
simulated impact, or elicited from a panel of practitioners with reported inter-rater agreement ---
would let the decomposition make an accuracy claim as well as an explanatory one.

\textbf{The Vulnerability dimension is the lightest of the four}, resting on reachability-style proxies
with no model of trust boundaries, privilege, or data sensitivity. A richer adversarial model would
strengthen the V attribution and broaden the framework's security relevance.

\textbf{Remediation is verified but not yet demonstrably effective.} The per-edit acceptance filter of
Section~\ref{sec:6.4} is implemented, which removes the possibility of an unverified regressing edit being applied
(Section~\ref{sec:6.7}). What it does not do is make the operators work: the risk reductions it achieves are real but
small (\(+0.0025\) to \(+0.0158\) SRI), and the operator set is narrow enough that its yield depends
heavily on whether the topology happens to contain a fan-out bottleneck. Expanding the operator set ---
along with deriving the acceptance multiplier \(\kappa\) from broader multi-seed variance data rather
than fixing it at 1.0 --- is needed before the prescriptive stage can claim practical value. One
specific gap follows: verification currently admits \emph{singletons}, each simulated on a graph
containing that edit alone, so verifying subsets rather than single edits is required before an
accepted policy can be called compositionally safe. That is engineering work rather than an open
research question, and we flag it as the immediate next step for this stage.

\textbf{The CI/CD gate is absolute, not yet delta-aware.} Section~\ref{sec:6.6} implements exit-code gating on a
candidate topology's full finding set; it does not yet diff that set against a merge-base topology,
so an architecture carrying an intentional, previously-accepted risk (a sole-source feed, a
deliberately unreplicated legacy broker) fails the build on every commit rather than only when a
change introduces something new. A delta-aware gate --- evaluate candidate and merge-base topologies,
block only on findings absent from the baseline --- together with a waiver register recording accepted
risk (entity, rule, expiry) so it stays auditable rather than silently re-suppressed on every run, is
the fix; neither is implemented today, and closing this gap is engineering work rather than an open
research question, similar in kind to the remediation gap above. A related, smaller gap is that
\texttt{detect\_antipatterns.py} itself requires a live Neo4j connection even though the underlying analysis
machinery is already usable through the database-free \texttt{MemoryRepository} --- wiring that path into the
packaged CLI is a prerequisite for running the gate without a database in an actual CI job.

\textbf{Edge-level ground truth is bounded by the cascade model.} Edge criticality is now measured by
removal rather than inferred from endpoints (Section~\ref{sec:8.2}), but the cascade routes no traffic over \texttt{RUNS\_ON}
or \texttt{CONNECTS\_TO} relations, so those edges measure as exactly zero regardless of their structural
role. Extending the cascade to express infrastructure-layer failure would close the same gap that
leaves Topic and physical Node components unlabelled.

\textbf{Relationship attribution is defined but not validated, and closing that gap is out of scope for
this submission.} Section~\ref{sec:4.7} gives D2 a measure with the same signature as D1, implemented and computed on
every derived dependency. What it does not have --- and, as the framework is currently constructed,
cannot have --- is the correlation-style evidence Section~\ref{sec:8.1} supplies for nodes: attribution is scored on
\texttt{DEPENDS\_ON} edges while the removal oracle severs raw structural edges, so the two are never defined
on a common population (Section~\ref{sec:4.7}). Re-simulating directly on the derived graph is not an available fix,
since the independence guarantee (Section~\ref{sec:5.3}) requires the simulator to operate only on
\(G_{\text{structural}}\). The one route that respects that guarantee is to track, for each derived
edge, which raw edges mediate it, and aggregate their measured impact onto it --- but the mediating
relations are many-to-many, so this requires a real modelling decision (how to aggregate) rather than
a mechanical lift, and we leave it for future work rather than attempt it here. Until it is done, the
relationship half of the framework rests on construction rather than on measurement, which we regard
as the most significant open gap in the diagnostic path.

\textbf{Finally, the endpoint for all of the above} is calibration against observed failure data from
instrumented deployments, which would convert this paper's comparative claims into absolute ones.

\hypertarget{conclusion}{%
\subsection{Conclusion}\label{conclusion}}

\label{sec:9.4}

We presented Software-as-a-Graph, a pre-deployment Static System Analysis (SSA) framework that
models distributed pub-sub middleware as a typed, weighted, directed multigraph and analyzes it
along two coupled axes: multi-dimensional quality attribution, which decomposes each component's
criticality into orthogonal, interpretable RMAV dimensions (integrating local code quality metrics),
and failure-impact analysis, which predicts cascade impact with both the interpretable composite and
a learned heterogeneous graph transformer, validated against discrete-event simulation under a
strict input--label independence guarantee. A prescriptive remediation stage generates topology-level
hardening edits from structure alone and verifies each one individually against the canonical
simulator, admitting it only when its benefit exceeds the simulator's own seed noise at every
propagation threshold; under that filter 162 of 332 candidates survive across the seven scenarios,
but the resulting risk reductions are small (\(+0.0025\) to \(+0.0158\) SRI) and concentrated in the two
topologies with pronounced fan-out structure, which we report as the substantive --- and qualified ---
result of that stage (Section~\ref{sec:6.4}, Section~\ref{sec:6.7}).

Integrated directly into pipelines as a delta-aware, blocking CI/CD Quality Gate, the framework
verifies architectural changes and blocks regression in seconds, bridging the ``Architecture-Code
Gap'' at commit time. Across a synthetic scenario suite, the framework establishes a scope condition
on where typed graph learning pays --- it leads a training-free QoS-weighted centrality on both ranking
(\(\rho = 0.608\) vs \(0.521\) out of distribution) and identifying \emph{which} components belong on a
shortlist (\(F_1@K = 0.465\) vs \(0.308\)), both measured only after repairing that baseline, which had
been computing no QoS weighting at all. Alongside that, measuring edge criticality by
removal rather than inferring it from endpoints shows most individual links to be replaceable and
exposes a class of relations the cascade model cannot express at all, and stratified rather than
pooled reporting caught a distortion that moved a headline figure by 0.38. By taking the \emph{type} of
every component and dependency seriously, the framework recovers structure that untyped,
single-dimensional methods discard, and does so at the point in the lifecycle where it is most
valuable: before the system runs.

\begin{center}\rule{0.5\linewidth}{0.5pt}\end{center}
